\subsection{Two actual age tags and cancellation of a common channel}
The retained source preparations at $T_1=T_0/4$ and $T_2=T_0/2$ encode
actual current ages $t_1=T_0/8$ and $t_2=T_0/4$.  Their increment is
$T_0/8$, rather than a full cycle.  Its phase is near $\pi/4$ and lies
strictly between zero and $\pi$ in every lane after the actual source disks
are paid.  This prefix prior supports unwrapping.  Repeated complex means
also require preparation of the same source relative phase against the
recorded reference; uncontrolled independent phase resets erase the mean.

Retain the original source convention $b_f=\int f a$, whose occupied
creation profile is $\overline f$.  A calibrated mode sorter transfers only
that canonical mode to the signal port and retains all complementary
outputs.  Known centering maps its two age-dependent profiles to the same
packet before the same complete annular channel.  For a common nonzero
complex coefficient $C$,
\[
 \beta_j=C\alpha_j,\qquad
 \beta_2\overline{\beta_1}
   =|C|^2\alpha_2\overline{\alpha_1}.
\]
The full channel phase cancels exactly.  This includes the actual annular
potential, rather than approximating it away.  Its previous $5.601$
microradian absolute free-comparator ceiling need not be paid as an
independent error in this strictly common experiment.  Channel magnitude
still determines brightness and noise.  If
$C_2/C_1=q e^{i\delta_{\rm ch}}$, $q>0$, the remaining phase bias is exactly
$\delta_{\rm ch}$.  A disk of relative radius $\rho<1$ about positive real
$q$ gives $|\delta_{\rm ch}|\leq\arcsin\rho$.  Comparing two physical
conditions requires their differential channel drift, not an assumption
that unrelated instrumental errors cancel.

A correctly matched real gate creates positive attenuation.  After the
specified mode processing, the temporal variance is
$\pi^2/(3\sigma^2)$.  With gate duration $G=2\pi/\sigma$ and center offset
$u$ as a fraction of its halfwidth,
\[
 \chi(u)\geq1-\frac{1}{3(1-|u|)^2},
\]
when this is positive.  The normalized positive-frequency projection of
the same windowed mode has real nonnegative overlap at least
$\sqrt\chi$.  A mismatched local-oscillator mode instead contributes its
actual complex overlap.  Mode leakage requires a number or energy bound:
$|\delta\alpha|\leq\delta_{\rm mode}\sqrt{N_{\rm leak}}$; a trace-distance
bound alone does not control an unbounded quadrature.  Hard leakage needs
its actual hard-state bound.  Buffering, frequency-shaping work, unused
outputs, source backaction and the late receiver reference remain paid
apparatus inputs.

\subsection{A nanoradian requirement, with its resource gap stated}
The heterodyne vacuum term remains in
$V\geq\langle N\rangle+1-|\langle a\rangle|^2$.  For mean error
$\epsilon$, blocks of $\lceil16V/\epsilon^2\rceil$ independent records and
coordinatewise median block means give the retained finite-variance bound.
Choose $\epsilon=39\,10^{-12}\alpha_-$ and 31 blocks.  All 24 complex means
(six lanes, two ages, two conditions) then have simultaneous failure below
\[
 24e^{-279/32}<0.003924.
\]
Four statistical phase disks, $0.020$ nrad differential channel mismatch,
$0.020$ nrad remaining apparatus/reference phase, and the actual intrinsic
source-field disks give a total below $0.196000001$ nrad.  The sufficient
fresh-preparation counts per complex mean, rounded upward, are:
\begin{center}\begin{tabular}{c|c|c}
 lane&preparation $T_0/4$&preparation $T_0/2$\\\hline
 01&$2.04632\,10^{59}$&$5.66920\,10^{58}$\\
 02&$4.76378\,10^{58}$&$1.31978\,10^{58}$\\
 03&$3.07426\,10^{56}$&$8.51703\,10^{55}$\\
 12&$1.77875\,10^{59}$&$4.92791\,10^{58}$\\
 13&$3.32718\,10^{56}$&$9.21774\,10^{55}$\\
 23&$3.63479\,10^{56}$&$1.00700\,10^{56}$
\end{tabular}\end{center}
These are explicit ideal-instrument requirements, not achieved precision,
a preparation rate or optimal copy lower bounds.  Both conditions require
their own samples.  In the actual $0<\omega<2\sigma$ band, the stated
combined reference/timing allocation requires
$\sum|\delta t|\leq0.020\,\mathrm{nrad}/(2\sigma)$, about $0.002026$ model
time units, unless a measured common timing error cancels.  The SI chart
remains a declared fixture, not a physical source-unit or detector
calibration.

The retained full-window negative phase margin is at least $4.30922$ nrad,
so this hypothetical $0.2$ nrad metrology budget fits it.  It does not assign
that full-window signal to the shorter tag increment.  That comparison
needs its perturbed source-emission bridge and actual mixed kernels.
The original full-window isolated susceptibility sign remains unresolved; the new signed prefix has a different, explicitly tiny time scale.  More samples
also do not remove the approximately 15 percent source-amplitude calibration
interval from distance inference or supply a geometric aperture law.

\subsection{Controlled amplitude, phase and instantaneous rate}
For an actual nonzero current $z_s$ write
\[
 \begin{gathered}
 y(t)=w(t)[1+u(t)],\qquad
 w(t)=A(t)e^{i\chi(t)}z_s(\tau(t)),\\
 A(t)>0,\qquad |u(t)|\leq\rho<1.
 \end{gathered}
\]
With clock phase $\theta=-\arg z$ and rate $\nu=\dot\theta$,
\[
 \begin{gathered}
 |\theta_y-\theta_s(\tau)+\chi|\leq\arcsin\rho,\\
 \nu_y=\dot\tau\,\nu_s(\tau)-\dot\chi
       -\operatorname{Im}\frac{\dot u}{1+u},\\
 |\nu_y-\dot\tau\,\nu_s+\dot\chi|
       \leq\frac{|\dot u|}{1-\rho}.
 \end{gathered}
\]
Positive amplitude alone changes neither phase nor rate.  A source-state
scalar Hamiltonian phase cancels from expectations altogether.  Physical
channel phase, time warp and additive residual instead need their actual
bounds.  The original instantaneous rate intervals can be propagated through
this formula; a rate sign requires the resulting interval to exclude zero.
An endpoint slowdown enclosure or a small phase disk supplies no bound on
$\dot u$ and does not certify an instantaneous slowdown sign.  Signal zeros
or loss of an amplitude floor invalidate phase tracking.

For the selected receiving form on the retained two-polarization radial
space, with $W=J^*NJ$, the exact one-particle operator is
$h_W=\sqrt\omega W\sqrt\omega$, and its source trace is
$v_{W,\alpha}=\sqrt\omega W F_\alpha$,
$F_\alpha=u_\alpha/\sqrt\omega$.  This trace is generally energy-dual,
not an ordinary Hilbert vector.  Use positive- and negative-half-energy
smearings.  Source-trace regularizations must converge in
$\|h_W^{-1/2}(v_n-v_W)\|$, with retained counterterms, convergent joint
forms, strong propagator convergence and uniform energy bounds.
This yields a weak matrix-element limit for original form-domain test
vectors; it does not require $a(F)$ to preserve that form domain.
For a creation-profile smearing $F$ the stationary weak relation is
\[
 \begin{aligned}
 a(F,T)={}&a(E_W(T)^*F,0)\\
 &+ig\sum_\alpha\int_0^T
 \langle F,E_W(T-s)v_{W,\alpha}\rangle S_\alpha(s)\,ds .
 \end{aligned}
\]
Choosing $F=E_W(T)F_{\rm soft}$ removes the initial hard homogeneous term
exactly for that adapted observable, while retaining all hard state and
radial hard/soft space.  Physical angular return belongs to the separate
canonical full-bundle problem.  It does not evaluate the remaining two-current kernels
or implement the altered detector.  Keeping the free kernel while merely
substituting a perturbed current is unjustified.  A nonautonomous extension also needs homogeneous half-energy propagation
bounds; shifted common-form existence alone does not supply them.
The common-form result above supplies the selected receiving dynamics; full canonical pair
creation remains a separate construction.

The actual real-time kernel is a complex spectral Fourier response.
It is not the positive Euclidean Laplace observable.  Physical relaxation
or cosmological $Q$ is phase-free only after proving the positive-multiplier
hypothesis with the required residual and time-warp bounds.  A convolution,
a changing canonical kinetic coefficient, or a changed source vertex can
add phase; no physical rate spectrum or $Q$ law is selected here.
